\documentclass{article}
\usepackage{graphicx} % Required for inserting images

\title{\vspace{-3em}TweetJoin (The This Game Crashes Game)}
\author{@DeltaClimbs \\\phantom{ }\\npub1c8nlcgd5l8qenesgdcy4vw0s794yujyy23z50n52v5ld0ddkcjnsppqsah}}
\date{March 2026}

\begin{document}

\maketitle

\noindent Due to the neurological tendencies of today, we provide the motivational inquiry as the 2nd section, though it was written in the reverse order.

\section{The TweetJoin Protocol (how to play)}

\begin{enumerate}
    \item \textit{Nothing is true, everything is permitted.} You may play by any rules you would like, this protocol is a mere possibility. A suggestion: either refrain from playing at all, or commit to playing strictly for at least 10 days to see what happens before leaving or perturbing.
    \item Form a private group chat anywhere of 3-8 friends, possibly with a dense amount of mutual ties; different things will happen if there is an oddball or in the group who might bridge to other clusters. If you want to go past 7-8 or so people, may make sense to split into 2 groups.
    \item Whenever you want to tweet, refrain from doing so directly; instead, post exactly what you would like to tweet (including any image or gif if desired, or specifying it as a reply or quote tweet to another) as a message within your group chat.
    \item Whenever you see someone post a message in the group, consider if it resonates, or perhaps even anti-resonates with you, and tweet it out, with quotation marks and attribution to the poster. In the attribution, would suggest not tagging with ampersand as the algorithm can easily detect and attack; could be practical to use perturbations of names, or even nicknames that are continuously rotated or shifted. See if you can have diversity/moderation as to who retweets.
    \item After tweeting, share in group so others know it has been posted and so author can then retweet.
    \item Consider every message fair game, exactly as written, with disregard to any metacommentary requesting it to be `outside the game' while within the group chat. Only consider metacommentary describing things such as those specifying that something is meant to be a retweet/reply.
    \item Refrain from retweeting anything that does not originate from a TweetJoin group -- if you run into people regularly you want to retweet, can you recruit them to a new TweetJoin group?   
\end{enumerate}

\section{Motivation and Inquiry}

Wherein there were once great promises and hopes of what might come of cyberspace, one might ask today, who was it who made such promises, and did they possess both strength and nobility to see these promises manifest? Or were they mere hopes, with little pieces of software operating within a broader framework within which human life and thinking itself had been encoded across many centuries? Might there be interstitial gaps within the scope through which the cypherpunk thinks -- that pesky, recurring problem, as to the nature of humans with their mathematical imperfections?\\ 

A great flattening of scope to that which was legible amidst the treacherous storm of combinatorial possibility; where might greater peaks be found, from within a system, which perhaps does not seek to give rise to nor elevate the highest men? -- a system of information controlled by machines of practicum, and this tasteless flattening I cannot abide: I say, let us \textit{Inject Human} where he is needed! Not for equality, but toward \textit{inequality,} and elevation of excellence amidst a substrate whose objective shall serve the purpose of each individual who commands within it, rather than tolerating enslavement to the substrate rules -- for are there not many men today who wish to be more than decorated slaves? Does this game teach the player something about a sly, roundabout way to insert capacity for authorship and creation? Let us create our own rules, irrespective of algorithms imposed... and shall we thereby discover an emerging type of man capable of transforming Leviathans to his will? \\

Doesn't everyone know something about Elon Musk he does not know about himself, and is there not something there that has leaked into the philosophy of twitter's algorithms? What does this philosophy breed?\\

Alas, what good will much further philosophizing serve us? Is it not the case that philosophy must be embodied? Can philosophy be interesting \textit{in-and-of-itself,} or must ideas connect to action, lest they fester in their impotence? I tire of all the theories, frameworks, and architectures, which are full of math and code, but devoid of life; this game is light on math, or determinateness, heavy on imposing players to take real risks in their real lives if they wish to play; what is really learned without that? Besides, do you really need me to pontificate, or can you guess what I might write if I were to go on for several more pages? Who says I am the best person to pontificate on this matter anyway? -- What might someone better versed in game theory or network science ponder? Does it really matter to attempt to explain a priori a game which no one is likely to care about since most games fail to continue anyway?\\

\noindent \textbf{Conjecture:} relational-first architectures overpower object-first architectures\\

\noindent \textbf{A matter of taste:} small group dynamics and emergence interest me more than large ensemble ``analysis'' (commentary)

\end{document}
